\chapter{Energy systems for a planet}
\label{vol12:ch02}

\epigraph{The energy transition is the largest engineering project in human history.}{Vaclav Smil, paraphrased from \emph{Energy and Civilization}, 2017}

\chapmeta{Half-life: medium-life. Archetypes: scaling, balance, ethics. Exercise target: 25. Project track: simulation (regional decarbonisation pathway). Hazard class: high (civilization-scale stakes; multi-decade lock-in). Reader path default: standard, with the projection sections core and the alternative-carriers and just-transition sections mastery. Named cases: Fukushima Daiichi (\cite{acc:tepco-fukushima-2012}); Texas February-2021 grid failure (\cite{acc:ferc-nerc-texas-2021}).}

\noindent
Energy is the master commodity of civilization. Every other infrastructure depends on it; every other engineered system either consumes it or produces it; every climate trajectory turns on it. The energy transition currently underway, from a roughly $80\%$ fossil-fuel global primary-energy mix to a target mix that brings net carbon emissions to zero by roughly mid-century, is the largest engineering project in human history. Its scale: roughly $600\,\text{EJ}$ of annual primary energy, several trillion dollars of annual capital investment, and a multi-decade institutional commitment across nearly every government and major industrial actor on Earth. Its difficulty: the existing system is itself the result of more than a century of optimisation against a different objective (cheap energy with externalised emissions). The transition replaces that system while keeping the lights on, the trains running, the steel mills operating, and the developing world's energy poverty receding. The engineering discipline this chapter describes is the discipline by which that transition is made, system by system, decade by decade.

\section{Global energy demand: history and projection}\pathtag{core}

The IEA's \emph{World Energy Outlook 2024} reports global primary energy consumption at roughly $630\,\text{EJ/yr}$ as of 2023, of which $80\%$ comes from fossil sources (coal, oil, natural gas). The trajectory since 1800 shows three transitions: from biomass to coal (roughly 1800-1900), from coal to oil-and-gas (roughly 1900-1980), and the partial overlay of nuclear and renewable additions (1980 onward). Each transition took roughly 50 years and did not displace the prior energy source so much as add to it; the absolute use of biomass and coal both continue to rise in some regions even as their fractional share of the global mix declines.

\begin{archetype}[Scaling: energy demand]
Global primary energy demand scales with population and per-capita energy use. Population: $\approx 8 \times 10^9$ as of 2024, projected to peak at $\approx 10 \times 10^9$ around 2080. Per-capita energy: $\approx 75\,\text{GJ/person/yr}$ globally, with developed-world figures of $150\,\text{GJ/person/yr}$ (Europe) to $300\,\text{GJ/person/yr}$ (USA, Canada, Australia) and developing-world figures of $20$-$50\,\text{GJ/person/yr}$. Convergence of developing-world per-capita energy to even half of the European level would roughly double global energy demand. The transition's task: deliver this demand without doubling emissions.
\end{archetype}

\begin{estimation}
\textbf{Question.} A country with $10^8$ population and current per-capita energy use of $50\,\text{GJ/yr}$ targets convergence to $120\,\text{GJ/person/yr}$ over 30 years while also targeting zero net emissions. Estimate the annual investment required, in fraction of GDP.

\textbf{Reasoning.} New annual energy demand: $10^8 \times 120 = 1.2 \times 10^{10}\,\text{GJ/yr} = 12\,\text{EJ/yr}$. Build-out over 30 years implies new capacity additions averaging $\approx 0.4\,\text{EJ/yr}$ per year over the 30 years. At a capital intensity of $\approx \$2$/W for clean generation and assuming a $30\%$ capacity factor (mixed renewables), the capital required per $\text{EJ}$ delivered is roughly $\$200$ billion. Annual investment: $\approx \$80$ billion. For a GDP of $\$1$ trillion (typical for a country at this energy-use level), this is $8\%$ of GDP. Conclusion: civilization-scale transitions at the planet's edge of the developing-developed transition require investment fractions that historically only wartime or major industrialisation episodes have sustained.
\end{estimation}

\section{The energy mix: where we are, where we go}\pathtag{core}

The energy mix the engineering profession inherits is not one mix but several. Final energy use divides into electricity ($\approx 20\%$ of final), industrial heat ($\approx 30\%$), transport fuel ($\approx 30\%$), buildings heat ($\approx 20\%$). The electricity sub-system is the furthest along in decarbonisation: in 2023, roughly $30\%$ of global electricity came from low-carbon sources (renewables $\approx 22\%$, nuclear $\approx 9\%$). The industrial-heat sub-system is the least along: roughly $90\%$ remains fossil, with steel, cement, ammonia, and process heat all dependent on combustion. The transport sub-system is partially electrifying (light-duty road vehicles) but largely not (heavy trucks, ships, aircraft). The buildings sub-system is regional: in some places (Norway, France for cold-climate heating) substantially electrified; in others (USA gas heating, China coal heating) substantially not.

The target mix at net zero by mid-century, per the IEA's net-zero scenario, allocates roughly $70\%$ of primary energy to electricity (a doubling from current); $15\%$ to bioenergy and synthetic fuels; $10\%$ to hydrogen; the remainder to small residual fossil with carbon capture. The implication: the global electricity system must grow by roughly $3\times$ in capacity and almost completely decarbonise its supply. This is the engineering problem of the century.

\begin{principle}[Sectoral specificity]
The energy transition is not one transition but several. The electricity sector's transition is roughly $30\%$ done; the industrial-heat sector's transition is roughly $5\%$ done; the transport sector's transition is roughly $10\%$ done. Each sector has its own technology, infrastructure, and institutional constraints. Engineering plans that treat the energy transition as a single problem miss the sectoral specificity that determines which technologies, which capital, and which timelines apply where.
\end{principle}

\section{Decarbonisation pathways}\pathtag{core}

A decarbonisation pathway is a time-resolved trajectory from current emissions to a target (typically net zero by 2050, or earlier for developed economies). Pathways are constructed in integrated assessment models (IAMs) such as MESSAGEix, REMIND, GCAM, IMAGE, which couple energy-system models to economy models to land-use models to climate models. The IPCC's AR6 scenarios database \cite{std:ipcc-ar6-wg3-2022} contains roughly 1,200 published pathways consistent with limiting warming to $1.5\degC$ or $2\degC$ by 2100.

Pathways share several robust features: rapid electrification of end uses (transport, buildings, light industry); aggressive build-out of low-carbon electricity (wind, solar, nuclear, hydro); decarbonisation of remaining fossil generation through carbon capture; demand-side efficiency improvements averaging $1$-$2\%/\text{yr}$; some role for bioenergy and direct air capture in offsetting residual emissions. Pathways differ in: the speed of fossil phase-out, the role of nuclear, the assumed availability of negative emissions, the assumed behavioural change, the assumed availability of carbon capture and storage.

\begin{archetype}[Balance: pathway choices]
Every pathway is a balance among speed, cost, risk, equity. A pathway that minimises cost may rely heavily on technologies (e.g., bioenergy with carbon capture and storage) whose feasibility at the required scale is uncertain. A pathway that minimises risk may require behavioural changes whose social feasibility is uncertain. A pathway that maximises speed may produce stranded assets and just-transition burdens that are politically destabilising. The engineering profession contributes to the analysis; the choice among pathways is a societal choice that engineering informs but does not determine.
\end{archetype}

\begin{estimation}
\textbf{Question.} A net-zero by 2050 pathway requires global low-carbon electricity capacity to reach roughly $30\,\text{TW}$ by 2050, from a 2023 starting point of $\approx 5\,\text{TW}$. Estimate the annual capacity additions required, and compare to recent historical additions.

\textbf{Reasoning.} Capacity gap: $30 - 5 = 25\,\text{TW}$, over 27 years (2023-2050). Annual additions: $\approx 0.9\,\text{TW/yr}$. Recent additions (2023): $\approx 0.45\,\text{TW/yr}$ globally. Conclusion: the pathway requires approximate doubling of current build rate, sustained for nearly three decades. Achievable in principle (solar PV additions have doubled every $\approx 5$ years for the past decade), but requires sustained supply chain, capital, and policy support.
\end{estimation}

\section{Storage at planetary scale}\pathtag{standard}

A grid powered by intermittent renewables (wind, solar) requires energy storage to balance generation and demand across timescales: sub-hourly (frequency regulation), daily (diurnal cycle), weekly (weather variability), seasonal (summer-winter mismatch). The storage technology landscape: lithium-ion batteries (sub-hourly to daily; roughly $\$150$/kWh as of 2024 and falling); pumped hydro (daily to weekly; site-limited but mature); compressed-air energy storage (daily; few sites operational); thermal storage (daily to weekly; emerging at industrial scale); hydrogen and synthetic fuels (weekly to seasonal; substantially less efficient round-trip but storable at very large scale). Each technology has a characteristic cost-per-kWh and cost-per-kW; together they form a portfolio that must match the discharge-duration distribution of the residual demand.

The seasonal storage challenge is the hardest. In northern Europe, winter electricity demand is roughly $50\%$ higher than summer demand, and solar output is roughly $25\%$ of summer values. A high-renewables grid must move energy from summer to winter, which requires storage capacities measured in tens of terawatt-hours. Lithium-ion is uneconomic at this duration; hydrogen and synthetic fuels are the leading candidates, with round-trip efficiencies of $30$-$40\%$ that imply large oversizing of generation.

\begin{archetype}[Scaling: storage duration]
Storage technologies span six orders of magnitude in discharge duration. Sub-hourly: capacitors, flywheels. Hourly: lithium-ion batteries. Daily-to-weekly: pumped hydro, flow batteries, thermal. Weekly-to-monthly: compressed-air, large thermal. Monthly-to-seasonal: hydrogen, synthetic fuels. Each technology occupies its niche by cost structure. Capital-cost-per-kW (power capability) and capital-cost-per-kWh (energy capability) differ across orders of magnitude. The grid's storage portfolio must match the residual-demand duration distribution, not just the peak.
\end{archetype}

\section{Grids: continental and intercontinental}\pathtag{standard}

The transmission grid is the energy system's circulatory infrastructure. The synchronous-AC grid was developed in the early 20th century and reached continental scale (Continental Europe, North America, China, India) by mid-century. High-voltage DC (HVDC) interconnects, developed since the 1950s, now allow point-to-point transfer over thousands of kilometres at multi-GW scale. A high-renewables grid amplifies the value of long-distance transmission: when the wind is not blowing in one region, it may be blowing 2,000 km away; when the sun sets in one time zone, it is still high in another.

The intercontinental supergrid concept (DESERTEC for Europe-Africa; Asian Super Grid for East Asia; various Atlantic and Pacific proposals) extends this logic to the planetary scale. Engineering feasibility is well-established; institutional and political feasibility is the binding constraint. The HVDC line from the Bakken Formation in North Dakota to Chicago, the $\pm 800\,\text{kV}$ Xiangjiaba-Shanghai line in China ($\approx 2{,}000\,\text{km}$, $\approx 6\,\text{GW}$), and the North Sea HVDC network all demonstrate the technical capability; the question is whether the political and institutional structures can support the cross-border investment, the operational coordination, and the conflict resolution required.

\begin{principle}[Transmission as substitute for storage]
Long-distance transmission and energy storage are partial substitutes. A renewable-heavy grid with strong transmission can absorb regional variability through geographic averaging; one without can do so only through storage. Optimal mixes typically use both. The engineering analysis is the cost trade-off between transmission build-out (a one-time capital investment with multi-decade life) and storage build-out (typically shorter-life and higher operating cost). The institutional analysis is harder: transmission crosses jurisdictional boundaries, storage usually does not.
\end{principle}

\section{Nuclear power and its institutional context}\pathtag{standard}

Nuclear power supplies roughly $9\%$ of global electricity, from approximately 440 operating reactors in 32 countries with total installed capacity of $\approx 390\,\text{GW}$ as of 2024. The technology has been mature since the 1970s; the issue is not principally technical but institutional. Capital costs for new builds in Western countries have risen from $\$2{,}000$/kW in the 1970s to $\$10{,}000$/kW for recent EPR and AP1000 projects, with construction times of 10+ years. In contrast, recent Chinese and South Korean builds have maintained $\$3{,}000$-$\$5{,}000$/kW with 5-7 year construction. The cost difference reflects regulatory regime, supply chain continuity, and learning-curve effects from sustained build programmes.

The Fukushima Daiichi accident of 2011 \cite{acc:tepco-fukushima-2012} is the most recent civilization-scale nuclear failure. The technical mechanism was a station blackout from tsunami inundation of backup diesel generators, leading to core damage in three reactors and offsite contamination across approximately $1{,}000\,\text{km}^2$. The institutional aftermath: Germany's accelerated phase-out, Japan's multi-year shutdown of its fleet, increased global regulatory caution. The engineering interpretation: the accident did not invalidate nuclear power but did invalidate any nuclear-power programme that did not take seriously the full range of beyond-design-basis events. The contemporary small modular reactor (SMR) designs, with passive cooling and lower power density, are partial responses to that lesson.

\begin{archetype}[Ethics: nuclear waste]
Nuclear power produces high-level waste with hazardous lifetimes of $10^4$-$10^5$ years. The engineering of geological repositories (Onkalo in Finland; Yucca Mountain proposed in USA; Bure proposed in France) is mature; the institutional structures for multi-millennial stewardship are not. The engineering profession's responsibility to future generations is at its sharpest here: the waste exists; the choice is between deep geological storage with high upfront engineering and indefinite surface storage with continuing institutional commitment. Either choice obligates institutions that will outlive the engineers who design them.
\end{archetype}

\section{Hydrogen, synthetic fuels, alternative carriers}\pathtag{mastery}

Hydrogen has been proposed as an energy carrier since the 1970s and re-proposed in successive waves. The current wave (since approximately 2020) is driven by the recognition that some sectors (steel, ammonia, long-distance shipping, long-haul aviation) cannot easily electrify and that hydrogen produced from low-carbon electricity (``green hydrogen'') offers a decarbonisation pathway. The economic challenge: green hydrogen costs in 2024 are roughly $\$4$-$\$8$/kg, against fossil-derived (``grey'') hydrogen at $\$1$-$\$2$/kg. Closing the gap requires cheap renewable electricity ($<\$30$/MWh), cheap electrolysers ($<\$300$/kW), and high capacity factors, all of which are achievable in principle but not yet at scale.

Synthetic fuels (Fischer-Tropsch from captured CO$_2$ and green hydrogen; methanol; ammonia; sustainable aviation fuel) are similarly placed. Their thermodynamic floor is high: electricity-to-liquid-fuel round-trip efficiency is roughly $30$-$50\%$ depending on the fuel and the use case. Their advantage is compatibility with existing infrastructure and engines. The engineering judgment that the profession is asked to make: which sectors should be served by direct electrification, which by hydrogen, which by synthetic fuels. The judgment depends on the technology learning curves over the next decade, which are uncertain.

\section{Failure: stranded-asset and just-transition issues}\pathtag{core}

The energy transition is a planned disinvestment from a multi-trillion-dollar fossil-fuel capital stock. Coal-fired power plants, gas pipelines, oil refineries, internal-combustion engine factories, fossil-fuel service stations, and the human capital trained to operate them all face economic obsolescence over the transition's timeframe. The stranded-asset literature \cite{paper:mccollum-stranded-2018} estimates global fossil-fuel asset stranding in a $2\degC$ pathway at $\$1$-$\$4$ trillion, concentrated in upstream oil and gas, coal, and associated infrastructure. The just-transition literature \cite{paper:newell-mulvaney-2013} adds the workforce displacement: hundreds of thousands of workers in coal mining, oil and gas production, refineries, and conventional auto manufacturing whose communities are economically dependent on the existing energy system.

The failure mode is the same as that documented in Volume X Chapter 6: technical-economic transitions that proceed without adequate institutional support for the affected workers and communities produce political backlash, transition reversal, and sometimes regulatory capture by the displaced incumbent. The Appalachian coal-region politics of the 2010s, the French ``yellow vest'' protests against fuel-tax increases, the German coal-region politics around the Hambach Forest dispute are all instances. The engineering profession's role: surface the distributive consequences in the transition planning, identify the engineering investments that retain employment value in transition regions (offshore wind in Aberdeen; geothermal in oil-and-gas regions; battery and electrolyser manufacturing in former auto towns), and engage with the institutional design that turns the engineering plan into a politically sustainable transition.

\begin{principle}[Just transition as engineering scope]
A decarbonisation pathway that ignores its just-transition consequences is not a complete engineering plan. The engineering profession's responsibility includes the assessment of distributive consequences and the design of transition pathways that retain employment and economic value in affected regions. The plan whose technical decarbonisation is achievable but whose political execution fails has failed at the engineering scope appropriate to civilization scale.
\end{principle}

\section{Project}\pathtag{core}

\begin{project}[Regional decarbonisation pathway model]
Choose a region with available energy data (a country, a US state, or a European country). Build a quantitative decarbonisation pathway from the region's current energy mix to net-zero by 2050. Specify the assumptions (technology costs, learning rates, behavioural change, available carbon sequestration). Compare to a published plan for the same region (IEA, national strategy, IPCC scenario). Identify the major points of disagreement and the engineering judgments that drive them. Estimate the capital investment required, the workforce implications, the distributive consequences. Document the just-transition strategy for the affected regions and workforces.
\end{project}

\section{Exercises}

\subsection*{Calculation}\pathtag{core}

\begin{exercise}
A country's electricity demand is $500\,\text{TWh/yr}$ with $60\%$ from natural gas at $0.4\,\text{tCO}_2/\text{MWh}$. Compute the annual CO$_2$ emissions and the fraction of a $2\degC$-consistent carbon budget allocated to the electricity sector this represents.
\end{exercise}

\begin{exercise}
A solar PV farm has nameplate capacity $1\,\text{GW}$ and capacity factor $25\%$. Compute the annual energy delivered. Compare to the output of a nuclear reactor of $1\,\text{GW}$ nameplate and $90\%$ capacity factor.
\end{exercise}

\begin{exercise}
An offshore wind farm has $500\,\text{MW}$ nameplate, $50\%$ capacity factor, and capital cost $\$3{,}500$/kW with $\$80$/kW/yr operating cost, $25$-year life, and discount rate $7\%$. Compute the levelised cost of electricity.
\end{exercise}

\begin{exercise}
A lithium-ion battery stores $100\,\text{MWh}$ with $90\%$ round-trip efficiency, $\$200$/kWh capital cost, $5{,}000$ cycle life, and $\$30$/MWh charging cost. Compute the levelised cost of storage per MWh delivered.
\end{exercise}

\begin{exercise}
A hydrogen electrolyser at $70\%$ efficiency uses $0.05$/kWh electricity. Compute the cost per kg of hydrogen produced (energy content of hydrogen: $33.3\,\text{kWh/kg}$).
\end{exercise}

\subsection*{Derivation}\pathtag{standard}

\begin{exercise}
Derive the relationship between renewable capacity factor, storage requirement, and curtailment fraction for a grid powered entirely by a single intermittent source.
\end{exercise}

\begin{exercise}
Derive the optimal mix of transmission and storage capacity for a two-region renewable system, assuming uncorrelated wind output between the regions and a given cost structure for both transmission and storage.
\end{exercise}

\begin{exercise}
Derive the carbon-budget allocation across the four major final-energy sectors (electricity, industrial heat, transport, buildings) under a $1.5\degC$ pathway, given baseline emissions and assumed sectoral decarbonisation rates.
\end{exercise}

\begin{exercise}
Derive the round-trip efficiency of a hydrogen storage system (electrolysis $\to$ compression/liquefaction $\to$ generation), and identify the dominant loss mechanism.
\end{exercise}

\subsection*{Estimation}\pathtag{standard}

\begin{exercise}
Estimate the global lithium demand for a net-zero pathway with $2$ billion electric vehicles by 2050. Compare to current production and known reserves.
\end{exercise}

\begin{exercise}
Estimate the land area required for global solar PV deployment under a net-zero pathway. Compare to current agricultural land and identify the conflict potential.
\end{exercise}

\begin{exercise}
Estimate the workforce displaced by coal phase-out in a major coal-mining region (Appalachia, Silesia, Shanxi). Identify the workforce categories most exposed.
\end{exercise}

\begin{exercise}
Estimate the steel demand for global wind turbine deployment under a net-zero pathway. Compare to current steel production and identify the steel-sector decarbonisation requirement.
\end{exercise}

\subsection*{Simulation}\pathtag{standard}

\begin{exercise}
Simulate one year of operation for a high-renewables grid with wind, solar, storage, and gas backup. Sweep the renewable-fraction parameter and identify the storage and backup requirements as a function of renewable penetration.
\end{exercise}

\begin{exercise}
Implement a learning-curve simulation for solar PV cost reduction with cumulative deployment. Sweep the learning rate and identify the policy implications for early-stage technology support.
\end{exercise}

\begin{exercise}
Simulate the temporal dynamics of a coal-to-renewables transition in a region with $20\,\text{GW}$ coal capacity, including stranded-asset writedowns, workforce displacement, and just-transition investment. Identify the policy interventions that minimise the political-economic backlash.
\end{exercise}

\subsection*{Design}\pathtag{standard}

\begin{exercise}
Design a $10\,\text{GW}$ wind-and-solar hybrid project with co-located battery storage and HVDC connection to a distant load centre. Specify the capacity ratio, the storage duration, the transmission rating, the capacity factor, the levelised cost.
\end{exercise}

\begin{exercise}
Design a workforce-transition programme for a coal-mining region of $50{,}000$ workers facing phase-out over 15 years. Specify the retraining, the new-industry attraction, the bridging income support, the institutional structure.
\end{exercise}

\begin{exercise}
Design an institutional structure for intercontinental supergrid governance. Specify the membership, the decision rules, the conflict resolution, the cost allocation, the regulatory harmonisation.
\end{exercise}

\subsection*{Judgement}\pathtag{mastery}

\begin{exercise}
Argue: the nuclear-power phase-out in Germany after Fukushima was an engineering and climate mistake. Argue the opposite. Identify the empirical evidence on each side; identify the value judgments that distinguish the positions.
\end{exercise}

\begin{exercise}
Discuss the role of carbon capture and storage (CCS) in net-zero pathways. Identify the engineering judgments about feasibility at the required scale; identify the moral hazard concerns about CCS as a license for continued fossil use; identify the implications for the engineering profession's recommendations.
\end{exercise}

\begin{exercise}
Discuss the just-transition challenge as an engineering responsibility. Identify the engineering decisions that can support just transitions; identify the institutional limits of what engineering can do; identify what the engineering profession should engage with beyond the technical scope.
\end{exercise}

\begin{exercise}
Argue: the energy transition's failure mode is political-economic, not technical; the engineering profession's contribution to its success requires explicit engagement with political-economic constraints. Argue the opposite. Use the cases of Germany's Energiewende, China's coal expansion, and the US Inflation Reduction Act as evidence.
\end{exercise}

\begin{exercise}
Discuss the engineering profession's response to nuclear power. Identify the technical, economic, and institutional arguments on each side; identify the engineering profession's responsibility to engage with the public deliberation about the technology's role in the energy transition.
\end{exercise}

\begin{exercise}
Discuss the Fukushima Daiichi accident as a beyond-design-basis nuclear-engineering failure. Identify the engineering decisions and the institutional context; identify the implications for the future of nuclear power in the global energy transition.
\end{exercise}
